% Part: first-order-logic
% Chapter: axiomatic-deduction
% Section: provability-propositional

\documentclass[../../../include/open-logic-section]{subfiles}

\begin{document}

\iftag{FOL}
      {\olfileid{fol}{axd}{ppr}}
      {\olfileid{pl}{axd}{ppr}}

\olsection{\usetoken{S}{derivability} en de logische tekens voor en, of en als--dan}

\begin{explain}
  Met de relatie~$\Proves$ schrijven we op wat met axiomatische
  deductie formeel uit aannamen volgt. We laten zien dat deze manier
  van afleiden een paar basisfeiten over de propositionele
  connectieven, dus de logische verbindingstekens, aankan. Zo gelden
  $!A \land !B \Proves !A$ en $!A, !A \lif !B \Proves !B$. Dat tweede
  patroon heet modus ponens. Deze feiten hebben we nodig voor het
  bewijs van de volledigheidsstelling.
\end{explain}

\begin{prop}\ollabel{prop:provability-land}
  \begin{enumerate}
  \item \ollabel{prop:provability-land-left} Er gelden twee
    gevolgtrekkingen: zowel $!A \land !B \Proves !A$ als
    $!A \land !B \Proves !B$.
  \item \ollabel{prop:provability-land-right} Omgekeerd geldt
    $!A, !B \Proves !A \land !B$.
  \end{enumerate}
\end{prop}

\begin{proof}
  \begin{enumerate}
    \item De eerste twee gevolgtrekkingen volgen met modus ponens uit
      respectievelijk \olref[prp]{ax:land1} en
      \olref[prp]{ax:land1}.
  \item Begin met \olref[prp]{ax:land3}. Pas modus ponens twee keer
    toe.
  \end{enumerate}
\end{proof}


\begin{prop}\ollabel{prop:provability-lor}
  \begin{enumerate}
  \item De formules $!A \lor !B, \lnot !A, \lnot !B$ zijn samen
    inconsistent.
  \item Er gelden twee gevolgtrekkingen: zowel
    $!A \Proves !A \lor !B$ als $!B \Proves !A \lor !B$.
  \end{enumerate}
\end{prop}

\begin{proof}
  \begin{enumerate}
  \item Uit \olref[prp]{ax:lnot1} krijgen we $\Proves \lnot !A \lif
    (!A \lif \lfalse)$ en $\Proves \lnot !B \lif (!B \lif
    \lfalse)$. Gebruik op beide de deductiestelling. Dan krijgen we
    $\{\lnot !A\} \Proves !A \lif \lfalse$ en $\{\lnot !B\} \Proves
    !B \lif \lfalse$. Gebruik vervolgens \olref[prp]{ax:lor3}. Dat
    geeft $\{\lnot !A, \lnot !B\} \Proves (!A \lor !B) \lif
    \lfalse$. Pas ten slotte de deductiestelling toe. Dan geldt
    $\{!A \lor !B, \lnot !A, \lnot !B\} \Proves \lfalse$.
  \item Gebruik \olref[prp]{ax:lor1} voor de eerste gevolgtrekking en
    \olref[prp]{ax:lor2} voor de tweede. Pas in beide gevallen modus
    ponens toe.
  \end{enumerate}
\end{proof}

\begin{prop}\ollabel{prop:provability-lif}
  \begin{enumerate}
  \item \ollabel{prop:provability-lif-left} Er geldt
    $!A, !A \lif !B \Proves !B$.
  \item \ollabel{prop:provability-lif-right}
    Er gelden bovendien twee gevolgtrekkingen: zowel
    $\lnot !A \Proves !A \lif !B$ als $!B \Proves !A \lif !B$.
  \end{enumerate}
\end{prop}

\begin{proof}
  \begin{enumerate}
  \item We kunnen dit meteen !!{derive}:
    \begin{derivation}
      1. & $!A$ & \Hyp\\
      2. & $!A \lif !B$ & \Hyp\\
      3. & $!B$ & 1, 2, \MP
    \end{derivation}
    
  \item Gebruik voor de eerste bewering \olref[prp]{ax:lnot2} en voor
    de tweede \olref[prp]{ax:lif1}. Gebruik daarbij ook de
    deductiestelling.
  \end{enumerate}
\end{proof}

\end{document}
